\chapter{NKCamera : lire une caméra}

Une caméra ressemble à un fichier vidéo qu'on lirait en direct. Elle n'en est pas
un, et la différence tient en une phrase : \textbf{avec un fichier, vous
choisissez ; avec une caméra, vous subissez}. Le format des pixels, la
résolution disponible, la mise au point, le sens de l'image --- tout cela est
décidé par le matériel et par le système, jamais par vous.

Ce chapitre explique donc moins « comment demander » que « comment recevoir ».

\section{La façade, en huit appels}

\begin{nkcode}{Kernel/Runtime/NKCamera/src/NKCamera/NkCameraSystem.h}
bool Init();
NkVector<NkCameraDevice> EnumerateDevices();
bool StartStreaming(const NkCameraConfig &config = {});
bool GetLastFrame(NkCameraFrame &outFrame);
void SetFrameCallback(NkFrameCallback cb);
void StopStreaming();
void Shutdown();
\end{nkcode}

\begin{nkretenir}
\textbf{Deux façons de recevoir les images, et elles ne se valent pas.}

\texttt{GetLastFrame} rend la \emph{dernière} image, celle qui est là au moment
où l'on demande. C'est ce qu'il faut pour un affichage : à 60 images par seconde
d'affichage et 30 de caméra, on redemande simplement deux fois la même.

\texttt{SetFrameCallback} vous appelle à chaque image, \textbf{depuis le fil de
la caméra}. C'est ce qu'il faut pour un enregistrement ou une analyse --- et ce
qu'il ne faut surtout pas employer pour dessiner.
\end{nkretenir}

\begin{nkpiege}
\textbf{Le rappel s'exécute sur un autre fil.} Y dessiner, c'est toucher le
contexte graphique depuis un fil qui ne le possède pas --- selon le backend, un
plantage, un écran noir, ou pire : cela marche chez vous.

C'est la même règle que le chapitre NKThreading : \emph{le calcul se
parallélise, le dessin non}. Le rappel copie ou analyse ; le fil principal
dessine.
\end{nkpiege}

\section{Le format n'est pas un choix}

C'est le point le plus important du chapitre, et il est écrit dans le code
lui-même :

% `escapeinside` : le panneau ne s'imprime pas dans la police du code, et la
% regle `literate` ne mord pas dans un bloc de code -- mesure faite le 02/09.
\begin{nkcode}[listing options={escapeinside={(*}{*)}}]{Kernel/Runtime/NKCamera/src/NKCamera/NkCameraTypes.h}
// (*\alerte{}*) TROIS CHAMPS RETIRES LE 2026-08-15 -- ne pas les reintroduire
// sans les cabler d'abord :
//
//   `outputFormat` -- 3 ecrivains, 0 lecteur. Trois applications
//     demandaient `NK_PIXEL_RGBA8` et recevaient du NV12 (Windows),
//     YUV420 (Android), BGRA8 (Cocoa/UIKit) ou YUYV (Linux) : le
//     format est IMPOSE par la plateforme, jamais choisi.
\end{nkcode}

\begin{nkretenir}
\textbf{Un paramètre qui n'est pas honoré est pire qu'un paramètre absent :
l'absence force à chercher, la présence dispense de vérifier.}

Trois applications écrivaient \texttt{outputFormat = NK\_PIXEL\_RGBA8}, en toute
bonne foi, et recevaient autre chose. Rien ne les avertissait --- le champ
existait, il acceptait la valeur, et il ne la lisait pas.

Le remède appliqué n'a pas été de faire fonctionner le champ (c'est impossible :
le format vient du pilote) mais de le \textbf{retirer}. Le format réellement
livré se lit sur \texttt{NkCameraFrame::format}, après coup.
\end{nkretenir}

\begin{nkcode}{Kernel/Runtime/NKCamera/src/NKCamera/NkCameraTypes.h}
// Un champ mort dans un en-tete public est une promesse que
// quelqu'un finira par croire.
\end{nkcode}

\begin{nkpiege}
Le même bloc en signale un autre, et celui-là n'a pas été retiré :

\begin{quote}
« \texttt{autoFocus} --- lu par le SEUL backend Android. Sur Windows, aucun code
de mise au point n'existe : le réglage y est ignoré en silence --- invisible sur
une webcam à focale fixe, ce qui explique que personne ne l'ait vu. »
\end{quote}

Retenez la cause de l'aveuglement : \textbf{le défaut était masqué par le
matériel de ceux qui auraient pu le voir.} Sur une webcam à focale fixe, un
réglage de mise au point qui ne fait rien produit exactement l'image attendue.
\end{nkpiege}

\section{L'image reçue, et ce qu'elle déclare}

\begin{nkcode}{Kernel/Runtime/NKCamera/src/NKCamera/NkCameraTypes.h}
struct NkCameraFrame {
        uint32 width = 0;
        uint32 height = 0;
        NkPixelFormat format = NkPixelFormat::NK_PIXEL_RGBA8;
        uint64 timestampUs = 0;
        uint32 frameIndex = 0;
        uint32 stride = 0;
        NkColorRange range = NkColorRange::NK_COLOR_RANGE_UNKNOWN;
        bool flipHorizontal = false;
        // ...
};
\end{nkcode}

\begin{nkretenir}
\textbf{\texttt{range} vaut « inconnu » par défaut, et c'est délibéré.} Le
commentaire d'origine le dit : « Non renseignée = INCONNUE, et c'est voulu : le
défaut ne prétend rien. »

C'est l'inverse du réflexe habituel, qui est de choisir une valeur « raisonnable
». Une plage supposée à tort décale toutes les couleurs de quelques pour cent ---
assez pour être faux, trop peu pour être vu. Un \emph{inconnu} déclaré force
l'appelant à décider ; un défaut plausible le dispense de savoir qu'il y avait
une question.
\end{nkretenir}

\begin{nkretenir}
\textbf{Et la conversion choisit sa formule d'après la PLAGE, jamais d'après le
format.} C'est écrit dans le même commentaire, à propos de
\texttt{ConvertToRGBA8}. Deux images au même format peuvent avoir deux plages
différentes ; se fier au format donnerait la bonne réponse la plupart du temps,
ce qui est la pire des situations --- un défaut qui ne se reproduit que sur
certaines caméras.
\end{nkretenir}

\section{Le miroir : une question de géométrie, pas de goût}

\begin{nkcode}{Kernel/Runtime/NKCamera/src/NKCamera/NkCameraTypes.h}
// Faux par defaut : l'image brute d'un capteur est geometriquement VRAIE,
// et c'est la seule utilisable pour de l'AR, de la calibration ou de la
// mesure. On ne miroite que pour un affichage de SOI, ou l'habitude prime
// sur la geometrie -- et cela reste une demande explicite.
bool flipHorizontal = false;
\end{nkcode}

\begin{nkretenir}
Une caméra frontale affichée telle quelle donne une image « à l'envers » pour
l'utilisateur, qui s'attend à un miroir. Mais une image miroitée est
\emph{géométriquement fausse} : un marqueur de réalité augmentée s'y place du
mauvais côté, une calibration y rend des paramètres symétriques faux, une mesure
y est inversée.

D'où le choix : \textbf{la vérité par défaut, le confort sur demande}. Et le
miroir est reporté \emph{dans l'image} (\texttt{NkCameraFrame::flipHorizontal})
afin que la conversion sache ce qui a été demandé --- au lieu que chaque
consommateur devine.
\end{nkretenir}

\section{Ce qu'il faut faire, dans l'ordre}

\begin{center}
\begin{tabular}{@{}rll@{}}
\toprule
& \textbf{Geste} & \textbf{Pourquoi} \\
\midrule
1 & \texttt{Init()} & le système peut refuser --- vérifiez \\
2 & \texttt{EnumerateDevices()} & il n'y a pas toujours une caméra \\
3 & \texttt{StartStreaming(cfg)} & la résolution demandée n'est qu'un souhait \\
4 & lire \texttt{frame.format} & \emph{c'est là} qu'on apprend ce qu'on reçoit \\
5 & convertir en RGBA8 & le rendu ne connaît que ça \\
6 & téléverser dans une texture & comme pour une vidéo \\
7 & \texttt{StopStreaming()}, \texttt{Shutdown()} & dans l'ordre inverse \\
\bottomrule
\end{tabular}
\end{center}

\begin{nkpiege}
\textbf{L'étape 2 n'est pas facultative.} Un ordinateur de bureau sans webcam, un
téléphone dont l'utilisateur a refusé la permission, une caméra déjà employée par
une autre application : les trois sont normaux, et les trois donnent une liste
vide.

Une application qui suppose qu'il y a une caméra affiche un écran noir sans
explication --- et l'utilisateur croit que le programme est cassé.
\end{nkpiege}

\begin{nkretenir}
À partir de l'étape 5, vous êtes exactement dans le chapitre NKMedia : des pixels
RGBA, une texture qu'on remplace, un sprite qui l'affiche. \textbf{Le rendu ne
sait pas distinguer une caméra d'un fichier vidéo}, et c'est voulu.

Une seule chose diffère : la caméra n'a pas de fin. Il n'y a pas de dernière
image, donc pas de « lecture terminée » --- seulement un arrêt que vous décidez.
\end{nkretenir}

\begin{nkbilan}
Une caméra ne se configure pas, elle se \emph{consulte}. Le format des pixels est
imposé par la plateforme et se lit sur l'image reçue, jamais dans la demande ---
un champ qui prétendait le contraire a été retiré parce qu'il faisait croire à
trois applications qu'elles recevaient du RGBA8. La plage de couleur vaut
« inconnu » par défaut, et la conversion s'appuie sur elle plutôt que sur le
format. Le miroir est faux par défaut, parce que l'image brute est la seule
géométriquement vraie. Et à partir de la conversion, tout redevient identique à
la lecture d'une vidéo.
\end{nkbilan}

\begin{nkquestions}
\begin{enumerate}
\item Pourquoi ne peut-on pas choisir le format des pixels d'une caméra ?
\item Où lit-on le format réellement reçu ?
\item Quelle différence entre \texttt{GetLastFrame} et le rappel d'image, et
      lequel convient à l'affichage ?
\item Pourquoi \texttt{range} vaut-il « inconnu » plutôt qu'une valeur plausible ?
\item Pourquoi le miroir est-il désactivé par défaut ?
\end{enumerate}
\end{nkquestions}

\begin{nkpratique}
Écrivez un aperçu caméra sur \texttt{NkCanvasGuiApp} : énumérez les
périphériques, laissez choisir dans une liste, démarrez le flux, convertissez et
affichez.

\textbf{Trois exigences}, et ce sont elles qui font l'exercice :
\begin{enumerate}
\item afficher en clair le format \emph{réellement} reçu, pas celui demandé ;
\item afficher un message lisible quand la liste des périphériques est vide ---
      et le vérifier en débranchant la caméra ;
\item un bouton « miroir » qui bascule \texttt{flipHorizontal}.
\end{enumerate}
\end{nkpratique}

\begin{nkdemo}
Prouvez que le format demandé n'est pas le format reçu.

Demandez explicitement \texttt{NK\_PIXEL\_RGBA8} dans votre configuration ---
si le champ n'existe plus, notez-le et passez à la suite. Puis affichez
\texttt{frame.format} et \texttt{frame.stride} pour chaque image.

Comparez ensuite \texttt{stride} à $\texttt{width} \times 4$ : s'ils diffèrent,
vos pixels ne sont pas du RGBA8 contigu, et une copie ligne par ligne qui
ignorerait le pas produirait l'image inclinée du chapitre NKImage.

\textbf{Ce que la démonstration doit établir} : que l'écart existe sur
\emph{votre} machine, avec un chiffre. « Le format peut différer » est une
lecture ; « j'ai demandé RGBA8 et reçu NV12, stride 640 pour width 640 » est une
mesure.
\end{nkdemo}

\begin{nklogique}
Votre aperçu caméra fonctionne sur votre ordinateur portable et donne une image
verdâtre sur celui d'un collègue.

\begin{enumerate}
\item Énumérez au moins quatre causes possibles, en distinguant celles qui
      viennent du format, de la plage de couleur, et de la conversion.
\item Le fait que l'image soit \emph{verdâtre} plutôt que noire ou bruitée
      élimine-t-il quelque chose ? Justifiez.
\item Proposez une mesure qui tranche --- et dites ce qu'elle affiche dans chaque
      cas. Puis expliquez pourquoi « ça marche chez moi » est ici l'information
      la plus trompeuse du problème, et non la plus rassurante.
\end{enumerate}
\end{nklogique}
